\maketitle
\makesignature

\ifproject
\begin{abstractTH}
โครงงานนี้นำเสนอการพัฒนาระบบวัดความสูงของต้นไม้โดยใช้เทคโนโลยี Ultra-Wideband (UWB) ร่วมกับแพลตฟอร์มโดรน 
โดยมีวัตถุประสงค์เพื่อศึกษาความเป็นไปได้ในการใช้ UWB เป็นทางเลือกแทนเทคโนโลยีที่มีต้นทุนสูง เช่น LiDAR 
และการประมวลผลภาพถ่ายทางอากาศ

ระบบที่พัฒนาขึ้นประกอบด้วยอุปกรณ์ UWB ประเภท Tag ที่ติดตั้งบนโดรน และ Anchor ที่ติดตั้งบนพื้นดิน 
โดยใช้หลักการ Time of Flight (ToF) ในการคำนวณระยะระหว่างอุปกรณ์ เพื่อนำไปประมาณความสูงของต้นไม้ 
ข้อมูลที่ได้จากการวัดจะถูกบันทึกและประมวลผลด้วยโปรแกรมที่พัฒนาด้วยภาษา Python

การทดลองดำเนินการในพื้นที่ป่าจริง โดยใช้ต้นไม้จำนวน 3 ต้นเป็นกรณีศึกษา 
ผลการทดลองถูกนำไปเปรียบเทียบกับข้อมูลอ้างอิงจาก 3D Point Cloud เพื่อประเมินความถูกต้องของระบบ 
โดยทำการวิเคราะห์ค่าความคลาดเคลื่อนและความสม่ำเสมอของการวัด

ผลการศึกษาแสดงให้เห็นว่า ระบบ UWB สามารถใช้วัดความสูงของต้นไม้ได้ในระดับที่น่าพอใจ 
อย่างไรก็ตาม ความแม่นยำของการวัดอาจได้รับผลกระทบจากสภาพแวดล้อม เช่น การบังของพุ่มไม้และความไม่เสถียรของสัญญาณ 
ซึ่งเป็นข้อจำกัดที่ควรพิจารณาในการนำไปใช้งานจริง
\end{abstractTH}

\begin{abstract}
This project presents the development of a tree height measurement system using Ultra-Wideband (UWB) technology integrated with a drone platform. 
The objective is to investigate the feasibility of using UWB as a cost-effective alternative to conventional methods such as LiDAR and photogrammetry.

The proposed system consists of a UWB tag mounted on a drone and an anchor placed on the ground. 
The distance between devices is calculated using the Time of Flight (ToF) principle and used to estimate tree height. 
Measurement data are collected and processed using a Python-based application.

Experiments were conducted in a real forest environment using three sample trees. 
The results were compared with reference data obtained from a 3D point cloud to evaluate system accuracy. 
Error and consistency analyses were performed to assess measurement reliability.

The results indicate that the UWB-based system can provide a reasonable estimation of tree height. 
However, measurement accuracy is affected by environmental factors such as signal obstruction from foliage and signal instability, 
which should be considered for practical deployment.
\end{abstract}

\iffalse
\begin{dedication}
This document is dedicated to all Chiang Mai University students.

Dedication page is optional.
\end{dedication}
\fi % \iffalse

\begin{acknowledgments}
ขอขอบพระคุณอาจารย์ที่ปรึกษาโครงงาน ผศ.ดร.กำพล วรดิษฐ์ และคณะกรรมการ รศ.ดร.ศันสนีย์ เอื้อพันธ์วิริยะกุล และ รศ.ดร.อัญญา อาภาวัชรุตม์ 
ที่ได้ให้คำแนะนำและข้อเสนอแนะอันเป็นประโยชน์ต่อการดำเนินโครงงานนี้ 

ขอขอบคุณผู้ให้คำแนะนำในการเลือกใช้อุปกรณ์ ESP32 UWB ซึ่งเป็นส่วนสำคัญในการพัฒนาระบบ

ขอขอบคุณผู้จัดทำชิ้นส่วน 3D Print สำหรับติดตั้งอุปกรณ์บนโดรน ซึ่งช่วยให้การติดตั้งมีความเหมาะสมและปลอดภัย

ขอขอบคุณผู้ให้คำแนะนำในการใช้งานระบบ WebODM และการประมวลผลข้อมูล 3D Point Cloud ซึ่งมีส่วนสำคัญต่อการวิเคราะห์ผลการทดลอง

สุดท้ายนี้ ขอขอบคุณภาควิชาวิศวกรรมคอมพิวเตอร์ ที่ได้สนับสนุนทรัพยากรและอุปกรณ์ที่จำเป็นสำหรับการศึกษาและการทดลอง
\acksign{2026}{3}{27}
\end{acknowledgments}%
\fi % \ifproject

\contentspage

\ifproject
\figurelistpage

\tablelistpage
\fi % \ifproject

% \abbrlist % this page is optional

% \symlist % this page is optional

% \preface % this section is optional
